\labsection{2}{ERP Benefits, Challenges, and Modular Integration}{sec:lab02}
\labcovers{ch:erp-development}

\begin{labproject}{Teach one ERP benefit and one ERP challenge using the system you just built}
\textbf{Coverage.} Tutorial 2 + Chapter 2.\\
\textbf{Core system.} Odoo 19 Community app structure.\\
\textbf{Goal.} Move from generic statements such as ``ERP improves efficiency'' to evidence-based explanations of how integration creates benefits and why implementation is difficult.

\courselabflow{Choose benefit}{Trace records}{Measure effect}{Explain challenge}



\begin{labmode}
Tutorial 2 uses Team Benefit and Team Challenge. Odoo is used here as a concrete reference architecture: each team must point to at least one app-to-app connection or setup choice that demonstrates the selected benefit/challenge.
\end{labmode}

\medskip\noindent\textbf{Part A --- Team Benefit.}

Choose one ERP benefit such as shared data, faster process flow, management visibility, reduced duplicate entry, better inventory accuracy, or integrated reporting. Explain:
\begin{enumerate}
  \item the old fragmented problem;
  \item the ERP mechanism that changes it;
  \item the Odoo record or connection that demonstrates the mechanism;
  \item one metric that could show the benefit.
\end{enumerate}

Example structure: ``A Sales Order creates a linked Delivery and can create an Invoice. This reduces repeated entry of product, quantity, customer, and price. Measure the percentage of invoices requiring manual correction.''

\courselabtrace{Turn a claim into inspectable evidence}{Claim}{Shared customer data reduces duplicate entry}{Mechanism}{Quotation and invoice refer to the same customer}{Measure}{Count corrections before and after the proposed process}



\textbf{Use the chapter while coding.} Revisit \hyperref[ch:erp-development]{the earlier chapter explanation}; match its operation to the variables and state changes below.

\textbf{Follow the supplied implementation.} Use those concrete records to explain integration. If a benefit is proposed rather than measured, name the metric and baseline needed to test it. Sharing a customer name alone does not prove the records are actually linked.
\courselabtrace{Trace one shared customer through the process}{1. Input / state}{Use the same customer from the CRM exercise rather than a second spelling or duplicate contact.}{2. Operation}{Follow its opportunity, quotation/order and downstream fulfilment records.}{3. Check}{Record the linked references at each handoff and identify the department responsible for the next action.}


\medskip\noindent\textbf{Part B --- Team Challenge.}

Choose one challenge such as data migration, training, change resistance, scope creep, poor process design, weak sponsorship, role/permission design, or excessive customisation. Explain:
\begin{enumerate}
  \item why it is difficult;
  \item which users/processes are affected;
  \item what could happen if it is ignored;
  \item how the implementation team should reduce the risk.
\end{enumerate}

\medskip\noindent\textbf{Part C --- Inspect modularity in Odoo.}

Open the app grid and identify which app owns each of these records: customer, opportunity, quotation, purchase order, receipt, BOM, manufacturing order, delivery, invoice, employee. Then draw arrows showing which records create or feed the next record.

\begin{keyidea}
The point is to see modularity without silos. The records live in different apps because the workflows differ, but they remain connected through shared master data and document relationships.
\end{keyidea}

\medskip\noindent\textbf{Part D --- Teach the other subgroup.}

Follow Tutorial 2: explain your selected topic, listen to the other subgroup, take notes, and write the required paragraph demonstrating what you learned.
\end{labproject}

\textbf{Chapter connection.} Revisit \hyperref[ch:erp-development]{the chapter explanation} and map its states to the steps below.
% Reused from chapters/02-erp-development.tex
\begin{labvisual}{Recall the chapter model}
\begin{labgraphic}
\begin{tikzpicture}[x=1cm,y=1cm]
\draw[draw=UTPMuted!45,line width=0.7pt] (-0.2,0) -- (15.2,0);
\foreach \x in {0.9,4.6,8.3,12.9}{\draw[draw=UTPMuted!45] (\x,-0.12)--(\x,0.12);}
\foreach \x/\era in {0.9/1960s, 4.6/1970s, 8.3/1980s, 12.9/1990s onward}{
  \node[procnote,below=1.5mm] at (\x,0) {\era};}
\node[procquiet,minimum width=3.1cm,minimum height=1.35cm,align=center] (a) at (0.9,1.35)
  {\bfseries Inventory control\\[-0.15em]\tiny What do we have\\\tiny in the store?};
\node[procstep,minimum width=3.1cm,minimum height=1.35cm] (b) at (4.6,1.35)
  {MRP\\[-0.15em]\tiny\mdseries Which materials,\\\tiny\mdseries and by when?};
\node[procstep,minimum width=3.1cm,minimum height=1.35cm,draw=UTPTeal,fill=UTPPaleTeal] (c) at (8.3,1.35)
  {MRP II\\[-0.15em]\tiny\mdseries Materials \emph{and}\\\tiny\mdseries capacity, cost, labour};
\node[procstep,minimum width=4.1cm,minimum height=1.35cm,draw=UTPGreen,fill=UTPPaleTeal] (d) at (12.9,1.35)
  {ERP\\[-0.15em]\tiny\mdseries Company-wide processes\\\tiny\mdseries on shared data};
\draw[procflow] (a)--(b); \draw[procflow] (b)--(c); \draw[procflow] (c)--(d);
\node[proccap,anchor=west,text width=7.0cm,align=left] at (-0.2,3.05)
  {\bfseries scope: one store room \quad$\longrightarrow$\quad one factory};
\node[proccap,anchor=east,text width=5.0cm,align=right] at (15.2,3.05)
  {\bfseries $\longrightarrow$\quad one enterprise};
\draw[draw=UTPBlue!35,line width=2pt,-{Stealth[length=3mm]}] (-0.2,2.6) -- (15.2,2.6);
\end{tikzpicture}
\end{labgraphic}\end{labvisual}



\begin{labquestion}{1}{Trace a shared customer}
Open the customer from the quotation in Lab 4. Return to the quotation, then follow its order and invoice links after Lab 5. Compare customer identity, product, quantity and price on each record.
\end{labquestion}
\begin{labanswer}\textbf{Inspect the result.} Record the document IDs and the matching fields. Distinguish a shared master record from a copied transactional value.\end{labanswer}

\begin{labquestion}{2}{Measure the proposed benefit}
Choose one measurable issue, such as correction count per ten quotations. Define the denominator and observation window before collecting a baseline. Repeat using the linked workflow.
\end{labquestion}
\begin{labanswer}\textbf{Inspect the result.} Show the two counts and the process change. If no baseline exists, label the metric as proposed rather than claiming a measured improvement.\end{labanswer}


\begin{verification}
\textbf{Explain and extend your result.} Choose one process delay and define a baseline measure, data source and proposed improvement. State what observation would support or contradict the claimed ERP benefit.


\begin{itemize}
  \item The chosen benefit names the old problem, the ERP mechanism, an Odoo record that shows it, and one metric.
  \item The chosen challenge names who it affects, what happens if it is ignored, and how to reduce the risk.
  \item Each of the ten listed records is assigned to the app that owns it, with arrows for the connections between them.
  \item The paragraph on the other subgroup's topic is written in your own words.
\end{itemize}
\end{verification}
